\documentclass[a4paper,11pt]{ltjsarticle}
\usepackage[margin=25mm]{geometry}
\usepackage{amsmath,amssymb,amsthm,mathtools}
\usepackage[unicode,colorlinks=true]{hyperref}
\newtheorem{theorem}{定理}[section]
\newtheorem{lemma}[theorem]{補題}
\newtheorem{proposition}[theorem]{命題}
\newtheorem{corollary}[theorem]{系}
\theoremstyle{remark}\newtheorem{remark}[theorem]{注意}
\renewcommand{\proofname}{証明}
\title{混合非可換cubicの特異境界におけるpoint scheme\\\large 六重点と有限平坦な衝突}
\author{}\date{2026年11月11日・v42}
\begin{document}\maketitle
\begin{abstract}
混合係数の判別式$\lambda^3+3=0$では元の可換cubicは特異になるが、
非可換中心商のpoint schemeの孤立部分は$\mathbb C[\lambda,h]$上
自由階数24の模型に延長できる。$h\ne0$の判別式上では三個の二重点と
三個の曲線的六重点が現れる。六重点の局所環は$\mathbb C[t]/(t^6)$、
$h=0$では$\mathbb C[t,z]/(t^3,z^2)$である。
滑らかなPoisson多様体の零点定理への無条件な拡張は主張しない。
\end{abstract}
% 2026-11-11 / v42
\section{混合係数の判別式上の孤立point scheme}
\label{pdb:sec}
混合量子化族の関係式は
\[
[y,z]=h(x^2+\lambda uy),\quad[z,x]=h(y^2+\lambda ux),
\quad[x,y]=hz^2,\quad [u,-]=[v,-]=0.
\]
中心三次元の可換化は
\[
\Omega_{h,\lambda}^{\rm ab}
=F_\lambda-\tfrac{2h^2}{3}xyz+\tfrac{h^2\lambda^2u^2}{6}z
 +\tfrac{h^4\lambda u}{18}z^2,\quad
F_\lambda=(x^3+y^3+z^3)/3+u^3+v^3+\lambda uxy.
\]
$h\ne0$で中心座標$u\ne0$のpoint moduleは定数列となるので、
$u=1$チャートの可換化の商環で表される。以下の環自体は$h=0$でも定義する。
\begin{equation}
\mathcal T=\mathbb C[\lambda,h,x,y,z,v]/
(x^2+\lambda y,\ y^2+\lambda x,\ z^2,\
1+v^3+\lambda xy/3+h^2z(\lambda^2/6-2xy/3)).
\label{pdb:eq:ring}
\end{equation}

\begin{theorem}[判別式を越える有限自由性]\label{pdb:thm:flat}
$\mathcal T$は$\mathbb C[\lambda,h]$上自由階数24で、基底は
\[
\{x^ay^bz^cv^d:0\le a,b,c\le1,\quad0\le d\le2\}.
\]
したがって$\lambda^3=-3$上でも有限平坦性は失われない。
\end{theorem}
\begin{proof}
$x^2+\lambda y,y^2+\lambda x$は総次数順序で
先頭単項式$x^2,y^2$を持つmonic Gr\"obner基底である。
唯一の$S$多項式$\lambda(y^3-x^3)$は
$x^3=y^3=-\lambda xy$により零へ還元される。
従って最初の二式の商は基底$1,x,y,xy$を持つ階数4の
$\mathbb C[\lambda,h]$自由加群。
$z^2$を加えて階数8、最後の$v$に関するmonic三次関係で
$1,v,v^2$を加えて階数24となる。この証明は分母を使わない。
\end{proof}

\begin{theorem}[三個の二重点と三個の六重点]\label{pdb:thm:sixfold}
$\lambda_0^3=-3$かつ$h_0\ne0$とする。
式\eqref{pdb:eq:ring}の$(\lambda,h)=(\lambda_0,h_0)$特殊化は
六つの支持点を持つ。
$x=y=0$上には$v^3=-1$の三個の長さ2の点があり、
各局所環は$\mathbb C[\varepsilon]/(\varepsilon^2)$である。
他の三点は
\[
 (x,y)=(-\lambda_0\zeta,-\lambda_0\zeta^2),\qquad \zeta^3=1
\]
に支持され、各局所環は$\mathbb C[t]/(t^6)$である。
従って長さは$3\cdot2+3\cdot6=24$。
同じ$\lambda_0$で$h=0$なら後者の局所環は
$\mathbb C[t,z]/(t^3,z^2)$となり、同型ではない。
\end{theorem}
\begin{proof}
$\lambda_0\ne0$なので最初の二式の解は
$(0,0)$または表示した三つの点である。
Jacobian行列式$4xy-\lambda_0^2$は前者で$-\lambda_0^2$、
後者で$3\lambda_0^2$だから四つの枝はetaleに分離する。
$x=y=0$では
$1+v^3+(h_0^2\lambda_0^2/6)z=0$、$z^2=0$であり、
$v^3=-1$の単純な三根それぞれで長さ2を得る。
他の枝では$xy=\lambda_0^2$、$\lambda_0^3=-3$なので
\[
 1+\lambda_0xy/3=0,\qquad
 \lambda_0^2/6-2xy/3=-\lambda_0^2/2.
\]
従って各局所環は
\[
 \mathbb C[v,z]/(z^2,v^3-cz),\qquad c=h_0^2\lambda_0^2/2.
\]
$c\ne0$なら$z=v^3/c$を消去し$\mathbb C[v]/(v^6)$、
$c=0$なら$\mathbb C[v,z]/(v^3,z^2)$。
両者のZariski接空間の次元はそれぞれ1,2である。
\end{proof}

\begin{corollary}[衝突の局所模型]\label{pdb:cor:double-cubic}
$\lambda^3+3=\delta$と置き、$\lambda\ne0$で
$(x,y)=(-\lambda\zeta,-\lambda\zeta^2)$の枝を固定する。
その有限部分の局所模型は
\[
 \mathbb C[\lambda,h,v,z]_{\rm loc}/
 (z^2,\ v^3+\delta/3-(h^2\lambda^2/2)z).
\]
$h\ne0$なら$z$を消去し
\[
\boxed{\mathbb C[\lambda,h,v]_{\rm loc}/((v^3+\delta/3)^2)}
\tag{DB}\label{pdb:eq:double-cubic}
\]
となる。すなわち$\delta\ne0$の三個の二重点が、
$\delta=0$で一個の六重点に合体する。
\end{corollary}
\begin{proof}
三つの非零枝では$xy=\lambda^2$なので、式
\eqref{pdb:eq:ring}の最後の方程式は
$v^3+\delta/3-(h^2\lambda^2/2)z=0$。
$h\lambda\ne0$なら
$z=2(v^3+\delta/3)/(h^2\lambda^2)$を代入し、$z^2=0$
から表示を得る。
\end{proof}

\begin{remark}[適用範囲]\label{pdb:rem:scope}
$\lambda^3=-3$では可換三次超曲面自身が特異なので、
滑らかなPoisson cubicの零点定理やBondal予想の
適用対象ではない。本節は中心座標非零チャートの
point moduleを記述する可換化環の平坦性と局所型を証明した。
$h\ne0$かつ$h^6\ne-27$では、別に得た
Sklyaninのpoint schemeの既知表示と併せて
全point scheme中の有限部分とみなせる。
$h=0$での環\eqref{pdb:eq:ring}は$B_{0,\lambda}$自身の
point scheme（可換cubic全体）とは異なる。
$h^6=-27$でのSklyanin成分の退化、任意のnc cubicへの拡張は未解決である。
\end{remark}

\end{document}
